\documentclass[11pt]{article}
\usepackage{amsmath,amsthm,amssymb}
\usepackage[margin=1in]{geometry}
\usepackage{hyperref}
\usepackage{enumitem}

\newtheorem{theorem}{Theorem}
\newtheorem{lemma}[theorem]{Lemma}
\newtheorem{corollary}[theorem]{Corollary}
\newtheorem{proposition}[theorem]{Proposition}
\theoremstyle{definition}
\newtheorem{remark}[theorem]{Remark}
\newtheorem{definition}[theorem]{Definition}

\newcommand{\Hq}{H_q}
\newcommand{\Sn}{S_n}
\newcommand{\Vlam}{V^{\lambda}}
\newcommand{\Om}{\Omega}
\newcommand{\hatl}{\widehat\lambda}
\newcommand{\elh}{\ell(\hatl)}
\newcommand{\Rp}{{R'}}
\newcommand{\SYT}{\mathrm{SYT}}
\newcommand{\tr}{\mathrm{tr}}

\title{Trace vanishing of $\Om^{(\lambda)}$ on $\Vlam$ from the multi-row meta-theorem}
\author{Clio}
\date{16 May 2026}

\begin{document}
\maketitle

\begin{abstract}
Let $\hatl$ be a partition with $\elh \ge 2$ and every row of length
$\ge 2$, and let $\lambda = (\hatl, 1^q) \vdash n$ with
$\tau(\hatl) := \max(0, q + 2\elh - 1 - |\hatl|) \ge 1$.  We prove
$\tr_{\Vlam}(\Om^{(\lambda)}) = 0$ as a direct corollary of the
2026-05-15 multi-row meta-theorem (Pillar~1: image-in-$E^-_1$),
combined with a one-line algebraic identity in the Iwahori--Hecke
algebra: $\Om^{(\lambda)}(T_1+1) = (q+1)\,\Om^{(\lambda)}$, which
follows from $(T_1+1)^2 = (q+1)(T_1+1)$ and the fact that
$\Om^{(\lambda)}$ ends in the rightmost factor $(T_1+1)$ of
$\Rp_n$.  Trace-vanishing is then a one-line corollary of the
proved meta-theorem.  We also note the stronger structural
corollary: in the Hoefsmit seminormal basis ordered by descent at
position~$1$, the matrix of $\Om^{(\lambda)}$ has block form
$\bigl(\begin{smallmatrix}0 & 0\\ X & 0\end{smallmatrix}\bigr)$,
in particular $\Om^{(\lambda)\,2} = 0$ on $\Vlam$.
\end{abstract}

\section{Setup}\label{sec:setup}

We follow the conventions of the May-15 multi-row meta-paper.
$\Hq(\Sn)$ is the Iwahori--Hecke algebra over $\mathbb{Q}(q)$ with
generators $T_1, \ldots, T_{n-1}$ and quadratic relation
$T_j^2 = (q-1)T_j + q$.  The eigenvalues of each $T_j$ on any module
are therefore $q$ and $-1$.  Set $S_j := T_j + 1$.  For $k \ge 1$,
\[
   \Rp_k \;:=\; S_{k-1} S_{k-2} \cdots S_1
   \;=\; (T_{k-1}+1)(T_{k-2}+1)\cdots(T_1+1) \quad (\Rp_1 := 1).
\]

For $\nu \vdash N$ with $\ell = \ell(\nu)$, set
\[
   \Om^{(\nu)} \;:=\; \Rp_{\ell+1}^{(\nu)} \Rp_{\ell+2}^{(\nu)} \cdots
   \Rp_N^{(\nu)}.
\]

Throughout this paper, $\hatl$ is a partition with $\elh \ge 2$ and
every row of length $\ge 2$.  We set
\[
   \lambda \;:=\; (\hatl_1, \ldots, \hatl_{\elh}, \underbrace{1, \ldots, 1}_{q})
   \;\vdash\; n,
   \qquad n = q + |\hatl|,
   \qquad \ell := \ell(\lambda) = \elh + q,
\]
and the threshold
\[
   \tau(\hatl) \;:=\; \max\bigl(0,\; q + 2\elh - 1 - |\hatl|\bigr).
\]

$\Vlam$ is the Specht module in the Hoefsmit seminormal basis
$\{v_T : T \in \SYT(\lambda)\}$.  For each $j \in \{1, \ldots, n-1\}$,
write
\[
   E_j^+ \;:=\; \ker(T_j - q)\!\mid_{\Vlam},
   \qquad
   E_j^- \;:=\; \ker(T_j + 1)\!\mid_{\Vlam}.
\]

\paragraph{Descent at position 1.}
For $T \in \SYT(\lambda)$, entry $1$ sits at cell $(1,1)$ and entry
$2$ sits at one of the cells $(1,2)$ or $(2,1)$.  We say $T$ is
\emph{descent-$0$ at position $1$} if $2$ is in row $1$ of $T$ and
\emph{descent-$1$ at position $1$} if $2$ is in column $1$ of $T$.
The axial distance between cells of $1$ and $2$ is then $+1$ or $-1$
respectively, so
\begin{equation}\label{eq:T1-action}
   T_1\,v_T \;=\;
   \begin{cases}
      q\, v_T, & T \text{ descent-}0\\
      -v_T, & T \text{ descent-}1.
   \end{cases}
\end{equation}
In particular $T_1$ acts diagonally on the Hoefsmit basis, and
\begin{equation}\label{eq:eigenspace-decomp}
   E_1^+\!\mid_{\Vlam} = \mathrm{span}\{v_T : T \text{ descent-}0\},
   \quad
   E_1^-\!\mid_{\Vlam} = \mathrm{span}\{v_T : T \text{ descent-}1\}.
\end{equation}

\section{Inputs}\label{sec:inputs}

\begin{theorem}[Pillar 1: multi-row meta-theorem]\label{thm:pillar1}
\textnormal{(\cite[Theorem 1]{multirowMay15}.)}
For $\hatl$ as above and $q \ge |\hatl| - 2\elh + 2$ (equivalently
$\tau(\hatl) \ge 1$),
\[
   \Om^{(\lambda)} \Vlam \;\subseteq\; \bigcap_{j=1}^{\tau(\hatl)} E_j^-\!\mid_{\Vlam}.
\]
In particular, $\Om^{(\lambda)} \Vlam \subseteq E_1^-\!\mid_{\Vlam}$.
\end{theorem}

\begin{lemma}[Right-absorption identity]\label{lem:right-absorb}
As an element of $\Hq(\Sn)$,
\[
   \Om^{(\lambda)}\,(T_1 + 1) \;=\; (q+1)\,\Om^{(\lambda)}.
\]
\end{lemma}

\begin{proof}
The Hecke quadratic relation $T_1^2 = (q-1)T_1 + q$ gives
\[
   (T_1+1)^2 \;=\; T_1^2 + 2T_1 + 1
              \;=\; (q+1)T_1 + (q+1)
              \;=\; (q+1)(T_1+1).
\]
The factor $\Rp_n = (T_{n-1}+1)(T_{n-2}+1)\cdots(T_1+1)$ ends with
$(T_1+1)$ on the right.  Therefore
\[
   \Rp_n \cdot (T_1+1)
   \;=\; (T_{n-1}+1)\cdots(T_2+1)\cdot(T_1+1)^2
   \;=\; (q+1)\,\Rp_n.
\]
Multiplying on the left by $\Rp_{\ell+1}\cdots \Rp_{n-1}$ and pulling
the scalar through gives $\Om^{(\lambda)}(T_1+1) = (q+1)\Om^{(\lambda)}$. $\qedhere$
\end{proof}

\begin{remark}
Lemma~\ref{lem:right-absorb} is an identity in $\Hq(\Sn)$, independent
of $\lambda$: it depends only on the rightmost factor of $\Om^{(\lambda)}$
being $(T_1+1)$.  In particular it implies $\Om^{(\lambda)} \cdot E_1^-
= 0$ on \emph{every} $\Hq$-module: for $v \in E_1^-$ we have
$(T_1+1)v = 0$, so $(q+1)\Om v = \Om(T_1+1)v = 0$, hence $\Om v = 0$
($q+1$ is a unit in $\mathbb{Q}(q)$).  This is the structurally
trivial right-vanishing on $E_1^-$, holding for every shape.
\end{remark}

\section{Main theorem}\label{sec:main}

\begin{theorem}[Trace vanishing on $\Vlam$]\label{thm:main}
Let $\hatl$ be a partition with $\elh \ge 2$ and every row of length
$\ge 2$, and let $\lambda = (\hatl, 1^q) \vdash n$ with
$\tau(\hatl) \ge 1$.  Then
\[
   \tr_{\Vlam}\!\bigl(\Om^{(\lambda)}\bigr) \;=\; 0.
\]
\end{theorem}

\begin{proof}
Compute the trace in the Hoefsmit seminormal basis:
\[
   \tr_{\Vlam}(\Om) \;=\; \sum_{T \in \SYT(\lambda)} [v_T]\bigl(\Om\,v_T\bigr),
\]
where $[v_T](w)$ denotes the coefficient of $v_T$ in the expansion of
$w$ in the seminormal basis.  We split the sum by descent at
position~$1$.

\emph{Case D0: $T$ descent-$0$.}  By
Theorem~\ref{thm:pillar1}, $\Om v_T \in E_1^-\!\mid_{\Vlam}$.  By
\eqref{eq:eigenspace-decomp}, $E_1^-\!\mid_{\Vlam}$ is spanned by
descent-$1$ basis vectors.  By linear independence of the Hoefsmit
basis, the coefficient of the descent-$0$ vector $v_T$ in any vector
of $E_1^-\!\mid_{\Vlam}$ is zero.  Hence $[v_T](\Om v_T) = 0$.

\emph{Case D1: $T$ descent-$1$.}  By \eqref{eq:T1-action},
$(T_1+1)v_T = 0$.  By Lemma~\ref{lem:right-absorb} as operators on
$\Vlam$,
\[
   0 \;=\; \Om(T_1+1)\,v_T \;=\; (q+1)\,\Om v_T.
\]
Since $q+1$ is a unit in $\mathbb{Q}(q)$, $\Om v_T = 0$.  In
particular $[v_T](\Om v_T) = 0$.

The trace vanishes term by term. $\qedhere$
\end{proof}

\begin{corollary}[Block-strict-lower form]\label{cor:block}
Order the Hoefsmit basis of $\Vlam$ as
$(v_T : T \text{ descent-}0\ ;\ v_T : T \text{ descent-}1)$, and write
$M(\Om)$ for the matrix of $\Om^{(\lambda)}$ in this basis.  Under
the hypotheses of Theorem~\ref{thm:main},
\[
   M(\Om) \;=\; \begin{pmatrix} 0 & 0 \\ X & 0 \end{pmatrix}
\]
for some matrix $X$.  Consequently $\Om^2 = 0$ on $\Vlam$,
$\mathrm{rank}(\Om^{(\lambda)}\!\mid_{\Vlam}) = \mathrm{rank}(X) \le \min(d_0, d_1)$
where $d_i = |\{T \in \SYT(\lambda) : T \text{ descent-}i\}|$, and
$\tr(\Om) = 0$.
\end{corollary}

\begin{proof}
For $T$ descent-$0$, $\Om v_T \in E_1^-\!\mid_{\Vlam}$ is supported only
on descent-$1$ basis vectors (Theorem~\ref{thm:pillar1} +
\eqref{eq:eigenspace-decomp}); column $T$ of $M(\Om)$ has zeros in
descent-$0$ rows.  For $T$ descent-$1$, $\Om v_T = 0$
(Lemma~\ref{lem:right-absorb} applied to $v_T \in E_1^-$); column $T$
is identically zero.  This gives the stated block form.  Then
$M(\Om)^2 = \bigl(\begin{smallmatrix}0&0\\X&0\end{smallmatrix}\bigr)
\bigl(\begin{smallmatrix}0&0\\X&0\end{smallmatrix}\bigr) = 0$, and
$\tr M(\Om) = 0 + 0 = 0$. $\qedhere$
\end{proof}

\section{Computational verification}\label{sec:verification}

The forward-direction trace vanishing was verified empirically (prior
to this proof) on six shapes via symbolic computation in
\texttt{sympy} with $q$ symbolic.  The two $\tau \ge 1$ shapes within
the multi-row meta scope are:
\begin{center}
\begin{tabular}{lccccl}
\hline
$\lambda$ & $\elh$ & $|\hatl|$ & $q$ & $\tau$ & $\tr_{\Vlam}(\Om^{(\lambda)})$\\
\hline
$(3,2,1,1,1)$    & $2$ & $5$ & $3$ & $1$ & $0$\\
$(3,2,1,1,1,1)$  & $2$ & $5$ & $4$ & $2$ & $0$\\
\hline
\end{tabular}
\end{center}

In addition, four $\tau = 0$ shapes ($(3,2,1)$, $(3,2,1,1)$, $(4,2,1)$,
$(3,3,1)$) were tested and have non-zero trace, consistent with the
forward direction (and providing evidence for the
still-open converse).

The script
\texttt{\textasciitilde/projects/scratch/2026-05-16-trace-vanishing-strategy/two\_sided\_signkill\_test.py}
also independently verifies the structural Corollary~\ref{cor:block}:
the matrix of $\Om^{(\lambda)}$ on $V^{(3,2,1,1,1)}$ in the Hoefsmit
basis has the block-strict-lower form, with rank $2$.

\section{Open: converse direction}\label{sec:open}

The converse direction
\[
   \tau(\hatl) = 0 \;\Longrightarrow\; \tr_{\Vlam}(\Om^{(\lambda)}) \ne 0
\]
remains open.  Empirical evidence: four $\tau = 0$ shapes verified
non-zero (above).  The proof in this paper does not address the
converse; it gives only the forward direction.  A natural next step
is leading-coefficient analysis of $\tr(\Om)$ as a polynomial in $q$,
or an explicit non-vanishing argument from the seminormal-basis
matrix at descent-$0$ diagonal entries.

\section{Methodological note}\label{sec:method}

This proof is a near-immediate corollary of the multi-row meta-theorem
plus the trivial right-absorption identity.  The discovery process
(2026-05-15 wake-2 through 2026-05-16 wake afternoon) is recorded in
the dream journal and need not concern us here, except for one
methodological observation: when a new conjecture appears, check
first whether it follows from existing pillars and an elementary
algebraic identity, before invoking heavy machinery (Markov traces,
JM products, categorical lifts).  In the present case, the
trace-vanishing conjecture admitted a one-paragraph proof from
already-proved Pillar~1 and a one-line lemma; the more elaborate
attack routes (q-binomial factorisation, JM-product involution,
Markov-trace cancellation) were not needed for the forward direction.

\begin{thebibliography}{9}
\bibitem{multirowMay15}
C.\ \emph{Extended sign-kill for general multi-row $\hatl$: the
meta-theorem at arbitrary $\ell(\hatl) \ge 2$.}
\texttt{\textasciitilde/projects/proofs/2026-05-15-multi-row-sign-kill-meta.tex}
(commit \texttt{7936707} on \texttt{clio-vega/proofs}, 2026-05-15.)
\end{thebibliography}

\end{document}
